\section{Data Logic Layer}

After presenting the main functionalities of the application, this chapter
focuses on the data logic layer, which is responsible for modeling and
managing the information handled by the system.

In particular, the structure of the database and the relationships between
entities are described through an Entity-Relationship schema.

\subsection{Entity-Relationship Schema}

The database is designed to represent the core concepts of the application
through a set of interconnected entities. The main entities are:
\texttt{User}, \texttt{Note}, \texttt{Course}, \texttt{Favorite}, and \texttt{Rating}.

The relationships between these entities can be summarized as follows:
\begin{itemize}
    \item A \texttt{User} can upload multiple \texttt{Note}s
    \item Each \texttt{Note} belongs to exactly one \texttt{Course}
    \item A \texttt{User} can mark multiple \texttt{Note}s as \texttt{Favorite}
    \item A \texttt{User} can assign a single \texttt{Rating} to a \texttt{Note}
    \item The \texttt{pair} (userId, noteId) is unique in the Rating entity.
\end{itemize}

\begin{figure}[H]
    \centering
    \fbox{\includegraphics[width=0.95\textwidth]{../diagrams/ERscheme.jpeg}}
    \caption{Entity-Relationship Schema of the Lecture Notes App}
    \label{fig:er-schema}
\end{figure}

The ER diagram in Figure~\ref{fig:er-schema} provides a visual representation
of the database structure and highlights the relationships between entities.

\subsection{Entity Description}

Each entity in the system stores specific information required to support
the application functionalities.

\textbf{Note.}
The \texttt{Note} entity contains the main information related to uploaded
lecture notes, including:
\texttt{id}, \texttt{authorId}, \texttt{authorUsername}, \texttt{courseId},
\texttt{courseName}, \texttt{title}, \texttt{description}, \texttt{uploadDate},
and \texttt{filePath}.
The \texttt{filePath} attribute represents the location of the PDF file,
which can be stored either on the server or on AWS S3 (as indicated by the
\texttt{software.amazon.awssdk:s3} dependency in the project configuration).

\textbf{User.}
The \texttt{User} entity stores authentication credentials and profile
information required for user identification and interaction with the system.

\textbf{Rating.}
The \texttt{Rating} entity represents the evaluation given by a user to a note.
It links a \texttt{User} to a \texttt{Note} through a numeric score.

\textbf{Favorite.}
The \texttt{Favorite} entity models the many-to-many relationship between
users and notes, allowing each user to maintain a personalized list of
favourite content.

\textbf{Course.}
The \texttt{Course} entity groups notes by subject, enabling better
organization and easier navigation within the application.